% TerraPulse AI — Mapa de cambio entre dos fechas NISAR (Chiapas y ROIs dibujadas en la web)
% Compilar con: pdflatex mapa_cambio.tex  (o latexmk -pdf), desde docs/latex/
\documentclass[11pt,a4paper]{article}
% Preámbulo común de la documentación LaTeX de TerraPulse AI.
% Se incluye con % Preámbulo común de la documentación LaTeX de TerraPulse AI.
% Se incluye con % Preámbulo común de la documentación LaTeX de TerraPulse AI.
% Se incluye con \input{preambulo} desde cada documento. Compilar con pdflatex (o latexmk -pdf).
\usepackage[utf8]{inputenc}
\usepackage[T1]{fontenc}
\usepackage{lmodern}
\usepackage[spanish,es-noshorthands,es-tabla,es-nodecimaldot]{babel}
\usepackage{amsmath,amssymb}
\usepackage[margin=2.3cm]{geometry}
\usepackage{array,booktabs,longtable}
\usepackage[normalem]{ulem}
\usepackage{xcolor}
\usepackage{enumitem}
\usepackage{listings}
\usepackage[hidelinks]{hyperref}

\setlength{\parindent}{0pt}
\setlength{\parskip}{0.5em}
\setlist{itemsep=2pt,topsep=2pt}
\renewcommand{\arraystretch}{1.25}
\Urlmuskip=0mu plus 1mu

% Símbolos frecuentes
\newcommand{\sz}{\ensuremath{\sigma^0}}
\newcommand{\gz}{\ensuremath{\gamma^0}}
\newcommand{\bz}{\ensuremath{\beta^0}}
\newcommand{\dsz}{\ensuremath{\Delta\sigma^0}}
\newcommand{\ok}{\ensuremath{\checkmark}}
\newcommand{\flecha}{\ensuremath{\rightarrow}}
\newcommand{\aprox}{\ensuremath{\approx}}
\newcommand{\por}{\ensuremath{\times}}

% Recuadro de nota (equivale a los ">" del Markdown)
\newenvironment{nota}
  {\begin{quote}\small\itshape}
  {\end{quote}}

% Bloques de código y comandos
\lstset{
  basicstyle=\ttfamily\small,
  breaklines=true,
  columns=fullflexible,
  keepspaces=true,
  frame=single,
  rulecolor=\color{black!30},
  backgroundcolor=\color{black!3},
  xleftmargin=4pt, xrightmargin=4pt,
  extendedchars=true,
  literate={á}{{\'a}}1 {é}{{\'e}}1 {í}{{\'i}}1 {ó}{{\'o}}1 {ú}{{\'u}}1
           {Á}{{\'A}}1 {É}{{\'E}}1 {Í}{{\'I}}1 {Ó}{{\'O}}1 {Ú}{{\'U}}1
           {ñ}{{\~n}}1 {Ñ}{{\~N}}1
}
 desde cada documento. Compilar con pdflatex (o latexmk -pdf).
\usepackage[utf8]{inputenc}
\usepackage[T1]{fontenc}
\usepackage{lmodern}
\usepackage[spanish,es-noshorthands,es-tabla,es-nodecimaldot]{babel}
\usepackage{amsmath,amssymb}
\usepackage[margin=2.3cm]{geometry}
\usepackage{array,booktabs,longtable}
\usepackage[normalem]{ulem}
\usepackage{xcolor}
\usepackage{enumitem}
\usepackage{listings}
\usepackage[hidelinks]{hyperref}

\setlength{\parindent}{0pt}
\setlength{\parskip}{0.5em}
\setlist{itemsep=2pt,topsep=2pt}
\renewcommand{\arraystretch}{1.25}
\Urlmuskip=0mu plus 1mu

% Símbolos frecuentes
\newcommand{\sz}{\ensuremath{\sigma^0}}
\newcommand{\gz}{\ensuremath{\gamma^0}}
\newcommand{\bz}{\ensuremath{\beta^0}}
\newcommand{\dsz}{\ensuremath{\Delta\sigma^0}}
\newcommand{\ok}{\ensuremath{\checkmark}}
\newcommand{\flecha}{\ensuremath{\rightarrow}}
\newcommand{\aprox}{\ensuremath{\approx}}
\newcommand{\por}{\ensuremath{\times}}

% Recuadro de nota (equivale a los ">" del Markdown)
\newenvironment{nota}
  {\begin{quote}\small\itshape}
  {\end{quote}}

% Bloques de código y comandos
\lstset{
  basicstyle=\ttfamily\small,
  breaklines=true,
  columns=fullflexible,
  keepspaces=true,
  frame=single,
  rulecolor=\color{black!30},
  backgroundcolor=\color{black!3},
  xleftmargin=4pt, xrightmargin=4pt,
  extendedchars=true,
  literate={á}{{\'a}}1 {é}{{\'e}}1 {í}{{\'i}}1 {ó}{{\'o}}1 {ú}{{\'u}}1
           {Á}{{\'A}}1 {É}{{\'E}}1 {Í}{{\'I}}1 {Ó}{{\'O}}1 {Ú}{{\'U}}1
           {ñ}{{\~n}}1 {Ñ}{{\~N}}1
}
 desde cada documento. Compilar con pdflatex (o latexmk -pdf).
\usepackage[utf8]{inputenc}
\usepackage[T1]{fontenc}
\usepackage{lmodern}
\usepackage[spanish,es-noshorthands,es-tabla,es-nodecimaldot]{babel}
\usepackage{amsmath,amssymb}
\usepackage[margin=2.3cm]{geometry}
\usepackage{array,booktabs,longtable}
\usepackage[normalem]{ulem}
\usepackage{xcolor}
\usepackage{enumitem}
\usepackage{listings}
\usepackage[hidelinks]{hyperref}

\setlength{\parindent}{0pt}
\setlength{\parskip}{0.5em}
\setlist{itemsep=2pt,topsep=2pt}
\renewcommand{\arraystretch}{1.25}
\Urlmuskip=0mu plus 1mu

% Símbolos frecuentes
\newcommand{\sz}{\ensuremath{\sigma^0}}
\newcommand{\gz}{\ensuremath{\gamma^0}}
\newcommand{\bz}{\ensuremath{\beta^0}}
\newcommand{\dsz}{\ensuremath{\Delta\sigma^0}}
\newcommand{\ok}{\ensuremath{\checkmark}}
\newcommand{\flecha}{\ensuremath{\rightarrow}}
\newcommand{\aprox}{\ensuremath{\approx}}
\newcommand{\por}{\ensuremath{\times}}

% Recuadro de nota (equivale a los ">" del Markdown)
\newenvironment{nota}
  {\begin{quote}\small\itshape}
  {\end{quote}}

% Bloques de código y comandos
\lstset{
  basicstyle=\ttfamily\small,
  breaklines=true,
  columns=fullflexible,
  keepspaces=true,
  frame=single,
  rulecolor=\color{black!30},
  backgroundcolor=\color{black!3},
  xleftmargin=4pt, xrightmargin=4pt,
  extendedchars=true,
  literate={á}{{\'a}}1 {é}{{\'e}}1 {í}{{\'i}}1 {ó}{{\'o}}1 {ú}{{\'u}}1
           {Á}{{\'A}}1 {É}{{\'E}}1 {Í}{{\'I}}1 {Ó}{{\'O}}1 {Ú}{{\'U}}1
           {ñ}{{\~n}}1 {Ñ}{{\~N}}1
}

\setlength{\parskip}{0.3em}

\title{TerraPulse AI\\[2pt]\large Mapa de cambio $\Delta\gamma^0$ entre dos fechas NISAR:\\ método, implementación y resultados}
\author{Documentación técnica del proyecto}
\date{18 de noviembre de 2026}

\begin{document}
\maketitle

\begin{abstract}
El mapa de cambio indica, píxel por píxel, cuántos decibeles cambió la retrodispersión del radar
entre dos fechas comparables de NISAR. Este documento describe el método (máscaras, multilook en
potencia lineal, conversión a dB, diferencia, z robusto y umbral), su implementación en
\path{pipeline/change.py} y \path{pipeline/roi_explore.py}, el servicio web que lo recalcula para
cualquier ROI dibujada y los resultados obtenidos con datos reales. Incluye el \emph{ruido de
referencia} de cada ROI, que indica qué tan confiable es un umbral fijo en dB en esa región.
\textbf{El resultado es exploratorio:} todavía no tiene validación independiente y no identifica
la causa física de ningún cambio.
\end{abstract}

\section{Qué es y qué no es}

\begin{itemize}
  \item \textbf{Es} una medición: $\Delta = \text{dB}_{t_2} - \text{dB}_{t_1}$ en cada píxel. Azul
        significa que en $t_2$ regresó menos señal; rojo, que regresó más; blanco, que no hubo cambio.
  \item \textbf{No es} un mapa de inundación, deforestación ni daño. Un cambio en dB describe un
        cambio en la \emph{respuesta del radar}. La misma subida puede deberse a humedad del suelo,
        rugosidad, vegetación o geometría. Distinguir la causa requiere fuentes auxiliares
        (paso~6 del plan de trabajo).
  \item \textbf{No es} exacto en un sentido demostrado. Sin etiquetas independientes no se reporta
        precisión, recall, F1 ni IoU; las cifras de área son \emph{conteos de píxeles sobre un
        umbral}, no áreas validadas.
\end{itemize}

\section{Datos y comparabilidad}

Se usan productos GCOV (nivel 2, retrodispersión geocodificada y corregida por terreno) de la
\emph{Frequency A} (banda L). Dos escenas solo se comparan si coinciden en producto, trayectoria,
marco, dirección de vuelo, polarización, resolución, CRID y madurez (\path{pipeline/compat.py}); no
se mezcla BETA con PROVISIONAL. Además, ambas fechas deben compartir la misma malla píxel a píxel
(CRS, transformación y forma). \textbf{Nunca se remuestrea para comparar}: si la malla difiere, el
cálculo se rechaza.

Las polarizaciones se leen de los metadatos de cada escena. En los casos de este documento son HH
y HV (adquisición \texttt{DHDH}).

\section{Cadena de procesamiento}
\label{sec:cadena}

Los pasos se aplican en este orden. Cada uno queda registrado en \path{change.json}.

\subsection{Máscaras}

Un píxel se descarta (NaN) si:
\begin{enumerate}
  \item la capa \texttt{mask} del producto vale 0 (no válido) o 255 (fuera de la adquisición);
  \item su centro cae fuera del polígono de la ROI (EPSG:4326 transformado a la proyección del
        producto, UTM).
\end{enumerate}
Todavía \textbf{no} se aplican máscaras de sombra, \emph{layover} ni pendiente con DEM.

\subsection{Magnitud}

GCOV entrega $\gamma^0$ en potencia lineal, normalizada por el área iluminada calculada con un
modelo de elevación. Por defecto se usa $\gamma^0$. Opcionalmente (solo en la línea de comandos)
se usa $\sigma^0 = \gamma^0 \cdot f$, con $f$ = \texttt{rtcGammaToSigmaFactor} del producto. Los
archivos dicen qué magnitud se usó.

\subsection{Multilook en potencia lineal}

La imagen SAR tiene \emph{speckle}: un ruido granular propio de la medición coherente. Aunque el
terreno no cambie, dos píxeles vecinos de la misma superficie pueden diferir varios dB. El
multilook $n\times n$ promedia bloques de píxeles \textbf{en potencia lineal}, nunca en dB:
\begin{equation}
  \bar{P}(I,J) = \frac{1}{|V_{IJ}|} \sum_{(i,j)\in V_{IJ}} P(i,j),
  \qquad
  |V_{IJ}| \ge 0.5\, n^2 ,
\end{equation}
donde $V_{IJ}$ son los píxeles válidos del bloque $(I,J)$. Si el bloque tiene menos de la mitad de
píxeles válidos, queda en NaN. Los bloques empiezan en la esquina superior izquierda del recorte y
las filas o columnas sobrantes del borde se descartan. La malla de salida es $n$ veces más gruesa:
con GCOV de 10\,m, $n=2$ da 20\,m, $n=5$ da 50\,m y $n=10$ da 100\,m.

Promediar dB sería incorrecto: el promedio de logaritmos no es el logaritmo del promedio y sesga el
resultado hacia valores bajos.

\subsection{Decibeles y diferencia}

\begin{align}
  \text{dB}_t(x,y) &= 10\,\log_{10} \bar{P}_t(x,y) \qquad (\bar{P}_t \le 0 \Rightarrow \text{NaN}),\\
  \Delta(x,y) &= \text{dB}_{t_2}(x,y) - \text{dB}_{t_1}(x,y).
  \label{eq:delta}
\end{align}
Si falta cualquiera de las dos fechas, $\Delta$ es NaN. La diferencia de dB equivale al logaritmo
de la razón de potencias: $\Delta = 10\log_{10}(\bar{P}_{t_2}/\bar{P}_{t_1})$.

\subsection{Puntaje z robusto}

Para describir cada píxel respecto al resto de la propia ROI se usa la mediana y la desviación
absoluta mediana (MAD), menos sensibles que la media y la desviación estándar a los píxeles que sí
cambiaron:
\begin{equation}
  s = 1.4826\cdot \operatorname{mediana}\bigl(|\Delta - \operatorname{mediana}(\Delta)|\bigr),
  \qquad
  z(x,y) = \frac{\Delta(x,y) - \operatorname{mediana}(\Delta)}{s}.
  \label{eq:z}
\end{equation}
El factor 1.4826 hace que $s$ coincida con la desviación estándar si $\Delta$ fuera normal.

\subsection{Clases por umbral fijo}

Con un umbral $u$ en dB:
\begin{equation}
  c(x,y) =
  \begin{cases}
    -1 & \Delta \le -u \quad\text{(baja la señal)}\\
    \phantom{-}0 & |\Delta| < u \quad\text{(estable)}\\
    +1 & \Delta \ge u \quad\text{(sube la señal)}
  \end{cases}
  \qquad (-128 \text{ sin dato}).
\end{equation}
El umbral solo decide qué píxeles entran en las cifras de área y en el GeoTIFF de clases. El mapa
de color muestra $\Delta$ continuo con la escala fija $\pm 6$\,dB, que no depende del umbral.

\subsection{Métricas}

Por polarización, sobre los píxeles válidos de la ROI: número de píxeles y área válida; media,
mediana, desviación estándar, MAD robusta y percentiles 2 y 98 de $\Delta$; área y porcentaje que
baja, sube o queda estable. El área de un píxel es $|a\cdot e|$ de la transformación afín (por
ejemplo, $50\times 50$\,m $= 0.0025$\,km$^2$).

\section{Umbral fijo y ruido de cada ROI}
\label{sec:ruido}

\subsection{El problema}

Con el mismo umbral en dB, una ROI más ruidosa marca más falsas alarmas. El ruido depende del
terreno (bosque, ciudad, agua, relieve), del par de fechas, del intervalo entre ellas y de la
resolución. En las ROIs dibujadas y guardadas al 18 de noviembre de 2026, el ruido $s$ del $\Delta$ HH
a 20\,m va de 1.09 a 1.53\,dB: el mismo umbral de 3\,dB equivale a entre 2.0 y 2.8 veces el ruido
según la ROI.

\subsection{Decisión}

Se probaron dos alternativas:
\begin{itemize}
  \item \textbf{Umbral relativo} ($|z|>3$): se adapta a cada ROI, pero el mismo criterio equivale a
        distintos dB en cada una, de modo que las cifras no se pueden comparar entre ROIs. Además, si
        un cambio real cubre gran parte de la ROI, infla $s$ y se diluye.
  \item \textbf{Umbral fijo en dB}: mide lo mismo en todas partes y es fácil de explicar, pero no
        dice cuándo está cerca del ruido.
\end{itemize}
Se eligió el \textbf{umbral fijo en dB como control} y se agrega, para cada ROI y polarización, un
\emph{ruido de referencia}:
\begin{align}
  \texttt{ruido\_dB} &= s \quad\text{(ecuación~\ref{eq:z})},\\
  \texttt{umbral\_sobre\_ruido} &= u / s,\\
  \texttt{umbral\_cerca\_del\_ruido} &= \bigl[\, u < 3\,s \,\bigr],\\
  \texttt{pct\_abs\_z\_mayor\_3} &= 100 \cdot \frac{\#\{|z|>3\}}{\#\{\text{válidos}\}}.
\end{align}
Si $u < 3s$, el manifiesto agrega una advertencia y la web muestra un aviso: parte de lo marcado
puede ser speckle, y conviene aumentar el multilook o el umbral. Estas cifras son una
\emph{referencia} de confiabilidad, no una exactitud.

\subsection{Efecto del multilook sobre el ruido}

Promediar $N$ muestras independientes reduce la varianza de la intensidad en un factor $N$. Los
píxeles de GCOV ya vienen promediados y están correlacionados entre vecinos, así que la reducción
real es menor que $1/\sqrt{N}$. En la ROI de prueba (tabla~\ref{tab:roi}), pasar de $2\times 2$ a
$5\times 5$ (6.25 veces más píxeles) bajó $s$ de 1.53 a 0.70\,dB, un factor de 2.2 frente al 2.5
que daría la raíz. Pasar a $10\times 10$ (25 veces) lo bajó a 0.39\,dB, un factor de 3.9 frente a 5. El
precio es la resolución: a 50\,m se pierden objetos de menos de unos 50\,m (caminos, casas
aisladas, ríos angostos).

\section{Implementación}

\subsection{Módulos}

\begin{longtable}{>{\raggedright}p{4.4cm} p{10.6cm}}
\toprule
Archivo & Función \\
\midrule
\endhead
\path{pipeline/preprocess.py} & Lee un recorte C1 (\path{data/raw/<scene_id>/}), aplica la máscara del producto y el polígono, multilook lineal y dB. \\
\path{pipeline/change.py} & $\Delta$, z robusto, clases, métricas, ruido de referencia, GeoTIFF y capas web. Entradas: recortes C1 (\texttt{run}), COG en dB (\texttt{change\_from\_db\_cogs}) o potencia lineal (\texttt{change\_from\_linear\_cogs}). \\
\path{pipeline/roi_explore.py} & Para una ROI dibujada: elige el par, llama a \texttt{etapa1} por fecha y calcula el cambio (\texttt{build\_change}); recalcula con otro multilook o umbral (\texttt{recompute\_change}) y limpia la caché. \\
\path{backend/scenes/roi_views.py} & Endpoints \texttt{/api/roi-jobs}; \texttt{RoiChange} recibe el recálculo. \\
\path{web/src/utils/changeLayer.ts} & Tipos del manifiesto, capa y métricas según la polarización elegida, resolución en metros y resumen de ruido. \\
\path{web/src/hooks/useRoiLayers.ts} & Carga la ROI y expone \texttt{recomputeChange}. \\
\path{web/src/components/ChangeControls.tsx} & Selectores de resolución y umbral, ruido de la ROI, $|z|>3$ y aviso. \\
\bottomrule
\end{longtable}

\subsection{Caso de estudio fijo (Chiapas)}

Desde la raíz del repositorio, con el par de \path{pipeline/config.yaml}:
\begin{lstlisting}
python -m pipeline.change                       # multilook 5, umbral 3 dB, gamma0
python -m pipeline.change --multilook 5 --umbral 3 --magnitud sigma0
\end{lstlisting}
Escribe \path{data/processed/cambio_<producto>_T<tray>_F<marco>_<CRID>_<t1>_<t2>_<magnitud>/}.

\subsection{ROIs dibujadas en la web}

\begin{enumerate}
  \item \texttt{POST /api/roi-jobs} prepara las capas de cada fecha (\texttt{etapa1}, multilook
        $2\times 2$ para inspección). \texttt{etapa1} conserva además la potencia lineal enmascarada
        a la resolución del producto en \path{layers/<granuleId>/geotiff/gamma0_<capa>.tif}.
  \item Si el par es comparable y la malla coincide, \texttt{build\_change} calcula el cambio desde
        esa potencia lineal con los valores por defecto: \textbf{multilook $5\times 5$ y umbral
        $\pm 3$\,dB}. El manifiesto queda en \texttt{result.json} $\rightarrow$ \texttt{change}.
  \item En la web, el botón \textbf{Cambio $\Delta$} muestra la capa de la polarización elegida (HH o
        HV). La barra de color incluye los controles de resolución (20 / 50 / 100\,m) y umbral
        ($\pm$1.5, 2, 2.5, 3, 4, 5, 6\,dB).
  \item Cada cambio de control llama a
\begin{lstlisting}
POST /api/roi-jobs/{id}/change   {"multilook": 10, "umbral_dB": 2}
\end{lstlisting}
        que recalcula \textbf{solo} el $\Delta$ a partir de la potencia lineal guardada, sin volver a
        ASF. Respuestas: 200 con el manifiesto; 400 si los valores no están permitidos o la ROI no
        tiene par comparable; 409 si la ROI no está lista o esa combinación se está calculando; 404
        si el trabajo no existe.
\end{enumerate}

Cada combinación queda en \path{layers/cambio/ml<n>_u<umbral>/} y se reutiliza. Se conservan la de
por defecto (\texttt{ml5\_u3}) y las cuatro más recientes. Un directorio \texttt{.lock} evita
calcular dos veces la misma combinación y se libera a los 5 minutos si quedó huérfano. Los valores
permitidos están acotados para limitar el costo de cada petición y el disco por ROI.

\subsection{Salidas y trazabilidad}

\begin{longtable}{>{\raggedright}p{4.4cm} p{10.6cm}}
\toprule
Archivo & Contenido \\
\midrule
\endhead
\path{delta_<pol>_dB.tif} & $\Delta$ en dB (float32, NaN sin dato), en la malla UTM tras el multilook. \\
\path{zscore_<pol>.tif} & $z$ robusto. \\
\path{cambio_<pol>.tif} & Clases $-1/0/+1$ ($-128$ sin dato). \\
\path{change.json} & Escenas, parámetros (magnitud, fórmula, máscara, multilook y método, tamaño de píxel, umbral, definición de z y del ruido), métricas por polarización, rutas, advertencias y versión del código. \\
\path{web/} & COG y PNG en EPSG:3857 (paleta azul--blanco--rojo, $\pm 6$\,dB) y \path{capas.json}. \\
\bottomrule
\end{longtable}
Los PNG son solo para ver. \textbf{Nunca se calculan diferencias sobre ellos}: los valores están en
los GeoTIFF y COG.

\section{Resultados con datos reales}

\subsection{Caso de estudio: Chiapas}

Par T62/F10 ascendente, P05023 PROVISIONAL, 2026-06-18 $\rightarrow$ 2026-10-16 (120 días), ROI de
\S4 de \path{AGENTS.md}. $\gamma^0$, multilook $5\times 5$ (50\,m), umbral $\pm 3$\,dB. Cifras de
\path{data/processed/cambio_GCOV_T062_F010_P05023_20260618_20261016_gamma0/change.json}, generado con
la versión \texttt{change-1} antes de agregar el ruido de referencia. La razón $u/s$ y el
porcentaje con $|z|>3$ se calcularon a partir de ese mismo archivo.

\begin{table}[h]
\centering
\begin{tabular}{lrr}
\toprule
 & HH & HV \\
\midrule
Área válida & 1469.2\,km$^2$ & 1469.2\,km$^2$ \\
Mediana de $\Delta$ & $+0.21$\,dB & $+0.24$\,dB \\
Ruido $s$ (1.4826\,MAD) & 0.91\,dB & 0.91\,dB \\
$u/s$ & 3.3 & 3.3 \\
Baja $\ge 3$\,dB & 6.73\,km$^2$ (0.46\,\%) & 3.33\,km$^2$ (0.23\,\%) \\
Sube $\ge 3$\,dB & 27.34\,km$^2$ (1.86\,\%) & 32.83\,km$^2$ (2.24\,\%) \\
Estable & 97.68\,\% & 97.54\,\% \\
$|z|>3$ & 2.71\,\% & 2.71\,\% \\
\bottomrule
\end{tabular}
\caption{Mapa de cambio de Chiapas (exploratorio).}
\end{table}

\textbf{Observación:} en ambas polarizaciones, el área que sube es varias veces mayor que la que
baja (unas 4 veces en HH y 10 en HV), y la mediana de $\Delta$ es positiva. El speckle tiende a
repartirse por igual hacia arriba y hacia abajo, así que esa asimetría es \emph{compatible} con un
cambio real de la respuesta radar, aunque no lo demuestra. El método
\textbf{no} permite decir a qué se debe. El intervalo de junio a octubre abarca la temporada de lluvias,
pero atribuirle el cambio sería una inferencia que requiere validación (por ejemplo, con GPM IMERG o
HLS).

\subsection{ROI dibujada: recálculo con distinta resolución}

ROI de 9.3\,km$^2$ cerca de Tuxtla Gutiérrez (trabajo \texttt{574924\dots}). Par T62/F10
ascendente, P05036, 2026-10-28 $\rightarrow$ 2026-11-09 (12 días), elegido por
\texttt{select\_scenes} como el par comparable más reciente.

\begin{table}[h]
\centering
\begin{tabular}{lrrrrrc}
\toprule
Resolución & Umbral & Pol. & Ruido $s$ & $u/s$ & Baja / sube & Aviso \\
\midrule
20\,m ($2\times2$)   & 3\,dB   & HH & 1.53\,dB & 2.0 & 2.96 / 2.88\,\% & sí \\
                    &         & HV & 1.61\,dB & 1.9 & 3.65 / 3.34\,\% & sí \\
50\,m ($5\times5$)   & 3\,dB   & HH & 0.70\,dB & 4.3 & 0.00 / 0.03\,\% & no \\
                    &         & HV & 0.74\,dB & 4.1 & 0.03 / 0.00\,\% & no \\
100\,m ($10\times10$) & 3\,dB  & HH & 0.39\,dB & 7.7 & 0.00 / 0.00\,\% & no \\
                    &         & HV & 0.40\,dB & 7.5 & 0.00 / 0.00\,\% & no \\
50\,m ($5\times5$)   & 1.5\,dB & HH & 0.70\,dB & 2.2 & 2.06 / 1.90\,\% & sí \\
                    &         & HV & 0.74\,dB & 2.0 & 2.94 / 2.30\,\% & sí \\
\bottomrule
\end{tabular}
\caption{Misma ROI y mismo par con distintas combinaciones de multilook y umbral.}
\label{tab:roi}
\end{table}

\textbf{Lectura:} a 20\,m con $\pm 3$\,dB se marca cerca del 3\,\% de la ROI en cada sentido, con
porcentajes casi iguales de subida y bajada. Esa es la huella del speckle, y el sistema lo señala
con el aviso. A 50\,m lo marcado prácticamente desaparece. Bajar el umbral a 1.5\,dB a 50\,m vuelve a
acercarlo al ruido, y el aviso reaparece. Cada recálculo tardó entre 0.1 y 1.5\,s, y una combinación
repetida se sirvió de la caché.

\textbf{Verificación:} a 20\,m, el $\Delta$ calculado desde la potencia lineal es idéntico bit a bit
al calculado antes desde los COG en dB de \texttt{etapa1} (misma malla, 23\,175 píxeles válidos,
diferencia máxima 0.0\,dB).

\section{Limitaciones y advertencias}

\begin{itemize}
  \item \textbf{Exploratorio:} sin validación independiente (paso~6 del plan). No se reporta exactitud.
  \item \textbf{Sin causa:} un $\Delta$ describe la respuesta del radar, no su origen físico.
  \item \textbf{Sin máscara topográfica:} no se excluyen sombra, \emph{layover} ni pendientes fuertes.
  \item \textbf{Madurez:} los productos PROVISIONAL están calibrados y parcialmente validados, no
        completamente validados.
  \item \textbf{Ruido de referencia:} $s$ se calcula sobre toda la ROI. Si un cambio real cubre una
        fracción grande de ella, $s$ crece y la referencia se vuelve conservadora.
  \item \textbf{Comparación entre ROIs:} las cifras en km$^2$ con el mismo umbral son comparables en
        dB, pero no en confiabilidad. Hay que leerlas junto con $s$ y $u/s$.
  \item \textbf{Escala de color:} fija en $\pm 6$\,dB; los valores fuera de ese rango se saturan en el
        PNG, pero no en los GeoTIFF.
\end{itemize}

\section{Pruebas}

\begin{itemize}
  \item \path{pipeline/tests/test_change.py}: diferencia y propagación de NaN, z robusto, clases,
        áreas, malla distinta, multilook en potencia lineal antes del dB, ruido de referencia y
        advertencia cuando el umbral está cerca del ruido (datos sintéticos).
  \item \path{pipeline/tests/test_roi_explore.py}: cambio por defecto a $5\times5$ y $\pm3$\,dB,
        recálculo con otra combinación, caché, validación de valores y limpieza de la caché.
  \item \path{pipeline/tests/test_roi_http.py}: endpoint de recálculo (200, 400, 404, 409) y URLs de
        los archivos.
  \item \path{web/tests/changeLayer.test.mjs}: capa y métricas por polarización, resolución en
        metros, opciones y resumen de ruido.
\end{itemize}
Las pruebas con datos sintéticos verifican el código, no la calidad del resultado sobre el terreno.

\end{document}
